\section*{Capítulo \ref{ch:g5:irreversible-changes} --- Cambios que no tienen vuelta atrás}

\begin{solution}{exo:g5:irreversible-changes:1}
Fundir chocolate: se puede deshacer (al enfriarse vuelve a solidificar).
Quemar una cerilla: no. \omterm{def:g2:mixing-and-dissolving:dissolve}{Disolver} sal: se puede deshacer (\omterm{def:g3:separating-mixtures:evaporate}{evaporando} el
agua). Tostar pan: no.
\end{solution}

\begin{solution}{exo:g5:irreversible-changes:2}
Un cambio en el que aparecen sustancias nuevas. En la cocina: cocinar
un huevo, hornear un bizcocho (también tostar pan, o una manzana
cortada que se pone marrón).
\end{solution}

\begin{solution}{exo:g5:irreversible-changes:3}
\omterm{def:g3:separating-mixtures:evaporate}{Evaporando} el agua: la sal se queda en el recipiente. \omterm{def:g2:mixing-and-dissolving:dissolve}{Disolver} es un
\omterm{def:g5:irreversible-changes:reversible}{cambio reversible}.
\end{solution}

\begin{solution}{exo:g5:irreversible-changes:4}
Un olor nuevo, y la aparición de grumos (un sólido) en la leche.
\end{solution}

\begin{solution}{exo:g5:irreversible-changes:5}
Los cambios de la izquierda: el azúcar que se \omterm{def:g2:mixing-and-dissolving:dissolve}{disuelve} en agua, el
hielo que se funde, la arena y la grava \omterm{def:g2:mixing-and-dissolving:mix}{mezcladas}. La doble flecha
significa que el cambio puede ir en los dos sentidos: se puede deshacer.
\end{solution}

\begin{solution}{exo:g5:irreversible-changes:6}
El aceite mantiene el \omterm{def:g4:air-a-mixture-of-gases:air}{aire} y el agua lejos del hierro de la cadena:
frena su oxidación.
\end{solution}

\begin{solution}{exo:g5:irreversible-changes:7}
La aparición de burbujas de gas en un líquido sin calentarlo.
\end{solution}

\begin{solution}{exo:g5:irreversible-changes:8}
Útiles: cocinar los alimentos, hornear el pan, quemar gas para la
calefacción. Perjudiciales: la oxidación del hierro, los alimentos que
se estropean, la madera que se pudre.
\end{solution}

\begin{solution}{exo:g5:irreversible-changes:9}
El frío hace más lentos los \omterm{def:g5:irreversible-changes:chemical-change}{cambios químicos} que agrían la leche, así
que se conserva más tiempo.
\end{solution}

\begin{solution}{exo:g5:irreversible-changes:10}
No. Las burbujas son vapor de agua, que vuelve a ser agua al enfriarse:
no se fabrica nada nuevo, y el cambio se puede deshacer. Las burbujas
solas no prueban que haya aparecido una sustancia nueva.
\end{solution}

\begin{solution}{exo:g5:irreversible-changes:11}
Una vela que \omterm{def:g4:what-a-fire-needs:burning}{arde} desprende calor y luz, y produce sustancias nuevas:
vapor de agua (un vaso frío se empaña encima de ella) y dióxido de
carbono. La cera que \omterm{def:g4:what-a-fire-needs:burning}{arde} no vuelve. En cambio, la cera derretida
vuelve a solidificarse al enfriarse.
\end{solution}
